\section{Scalar spectral measures}\label{sec:scalar-spectra}

The spectral arguments in Sections~\ref{sec:plane} and \ref{sec:pins}
use scalar covariance measures and bounded analytic matrix elements.
We construct these objects directly.  Inner products in this appendix
are conjugate linear in the first variable.  Haar measure on a torus
is normalized to have mass one.

\begin{lemma}[Commuting translations]\label{spectra:unitary}
Let $V_1,\ldots,V_d$ be commuting unitaries on a complex Hilbert space.
For vectors $a,b$ there is a unique finite complex measure $\mu_{a,b}$
on the $d$-torus such that
\[
 \int e^{i\mathbf n\cdot\theta}\,\mu_{a,b}(\dd\theta)
   =\langle a,V_1^{n_1}\cdots V_d^{n_d}b\rangle
 \qquad(\mathbf n\in\Z^d).
\]
The measure $\mu_{a,a}$ is positive, has mass $\|a\|^2$, and
$\|\mu_{a,b}\|_{\TV}\le\|a\|\|b\|$.  For Laurent polynomials $P,Q$,
\[
 \mu_{P(V)a,Q(V)b}
   =\overline{P(e^{i\theta})}Q(e^{i\theta})\mu_{a,b}.
\]
If $P_0$ is the orthogonal projection onto the common invariant
subspace, then $\mu_{a,b}(\{0\})=\langle P_0a,P_0b\rangle$.
\end{lemma}

\begin{proof}
For a vector $a$ and an integer $N\ge1$, consider the nonnegative density
\[
 K_{N,a}(\theta)=N^{-d}
 \left\|\sum_{0\le n_1,\ldots,n_d<N}
 e^{-i\mathbf n\cdot\theta}V_1^{n_1}\cdots V_d^{n_d}a\right\|^2.
\]
Its mass is $\|a\|^2$.  Its Fourier moment at $\mathbf m$ is
\[
 \prod_{r=1}^d(1-|m_r|/N)_+
       \langle a,V_1^{m_1}\cdots V_d^{m_d}a\rangle.
\]
Compactness of finite measures on the torus gives a subsequential
limit with the asserted moments.  Trigonometric polynomials are uniformly
dense in continuous functions, so those moments determine the limit
uniquely.  This constructs $\mu_{a,a}$ without an operator-valued measure.
Polarization constructs $\mu_{a,b}$ and gives its stated moments.

For every Borel set $A$ and any finite list of vectors, the matrix
$[\mu_{a_r,a_s}(A)]_{r,s}$ is positive semidefinite: its quadratic
forms are the positive measures of their linear combinations.
Thus $|\mu_{a,b}(A)|^2\le\mu_{a,a}(A)\mu_{b,b}(A)$.
Apply Cauchy--Schwarz to a finite measurable partition and take the
supremum over such partitions to obtain the total-variation bound.
The covariance identity for Laurent polynomials follows by checking
each Fourier moment, then using uniqueness.

Set
\[
 A_N=\prod_{r=1}^d\left(N^{-1}\sum_{k=0}^{N-1}V_r^k\right).
\]
These contractions act identically on common invariant vectors.
On a vector of the form $(I-V_r)b$, telescoping gives
$\|A_N(I-V_r)b\|\le2\|b\|/N$.
The orthogonal complement of the common invariant subspace is the
closed span of these ranges: a vector orthogonal to every range is
fixed by every $V_r^*$ and hence by every $V_r$.  Consequently $A_N$
converges strongly to $P_0$.  The polynomial covariance identity gives
\[
 \|A_Na\|^2
 =\int\prod_{r=1}^d\left|N^{-1}\sum_{k=0}^{N-1}e^{ik\theta_r}\right|^2
          \mu_{a,a}(\dd\theta).
\]
The integrand is bounded by one and converges to the indicator of the
origin.  Dominated convergence identifies its limit with
$\mu_{a,a}(\{0\})=\|P_0a\|^2$.  Polarization proves the cross identity.
\end{proof}

\begin{lemma}[A unitary and a positive contraction]\label{spectra:transfer}
Let $U$ be unitary and $S$ a positive self-adjoint contraction, with
$US=SU$.  For vectors $a,b$ there is a unique finite complex measure
$\gamma_{a,b}$ on the circle times $[0,1]$ satisfying
\[
 \int e^{im\theta}\lambda^j\,\gamma_{a,b}(\dd\theta,\dd\lambda)
    =\langle a,U^mS^jb\rangle
       \qquad(m\in\Z,\ j\ge0).
\]
The diagonal measures are positive, both endpoint atoms are retained,
and $\|\gamma_{a,b}\|_{\TV}\le\|a\|\|b\|$.
\end{lemma}

\begin{proof}
For $0\le k\le N$ put
\[
 B_{N,k}=\binom Nk S^k(I-S)^{N-k}.
\]
These operators are positive and sum to the identity.  To check
positivity without measurable functional calculus, first note that
$S$ and $I-S$ are positive contractions.  The scalar binomial series
for $\sqrt{1-t}$ has absolutely summable coefficients on $|t|\le1$.
It therefore gives norm-convergent self-adjoint square roots of $S$
and $I-S$, by substituting $I-S$ and $S$, respectively.  They commute,
being limits of polynomials in $S$.  Their suitable powers exhibit
each $B_{N,k}$ as a square of a self-adjoint operator.  The binomial
theorem gives $\sum_k B_{N,k}=I$.

The same square-root series applies to $B_{N,k}/\|B_{N,k}\|$ when
this operator is nonzero.  It gives $B_{N,k}^{1/2}$, which commutes
with $U$ because all its approximating polynomials do.  The zero
operator has square root zero.

For each $k$ apply the one-dimensional positive kernels in the proof
of Lemma~\ref{spectra:unitary} to $B_{N,k}^{1/2}a$, using $U$ and
the same integer $N$.  Place the resulting circle measure at
$\lambda=k/N$ and sum over $k$.  The obtained positive measure has
mass $\|a\|^2$, and its $(m,j)$ moment is
\[
 (1-|m|/N)_+
 \left\langle a,U^m\sum_{k=0}^N(k/N)^jB_{N,k}a\right\rangle.
\]
For every fixed $j$ the operator sum converges in norm to $S^j$.
Indeed express $k^j$ as a fixed linear combination of falling
factorials $(k)_r$; the binomial identity
$\sum_k(k)_rB_{N,k}=(N)_rS^r$ shows that the coefficient of $S^j$
tends to one and all lower coefficients tend to zero.
Compact extraction now yields the desired positive diagonal measure.
Products of trigonometric polynomials in $\theta$ and ordinary
polynomials in $\lambda$ approximate continuous functions uniformly,
so the moments uniquely determine it.  Polarization and the
positive-matrix argument from Lemma~\ref{spectra:unitary} give the
cross measures and their total-variation bounds.  The construction
takes place on the compact interval $[0,1]$; neither endpoint is removed.
\end{proof}

\begin{corollary}[Bounded analytic matrix elements]\label{spectra:analytic}
For the positive contraction $S$ above and vectors $a,b$, there is a
bounded holomorphic function $F_{a,b}$ on $\Re w>0$ such that
\[
 F_{a,b}(j)=\langle a,S^jb\rangle\quad(j=1,2,\ldots),
 \qquad |F_{a,b}(w)|\le\|a\|\|b\|.
\]
It is obtained by integrating $\lambda^w$ against the $[0,1]$
marginal of $\gamma_{a,b}$, with value zero at $\lambda=0$.
\end{corollary}

\begin{proof}
Take $U=I$ in Lemma~\ref{spectra:transfer} if no tangential operator
is present.  For $\lambda>0$ put
$\lambda^w=\exp(w\log\lambda)$ and set it equal to zero at zero.
When $\Re w>0$ its absolute value is at most one, giving the bound
and the positive-integer identities.  On each closed smaller
half-plane $\Re w\ge\eta>0$, every
$\lambda^\eta|\log\lambda|^r$ is bounded for integer $r\ge0$.
Dominated differentiation against the finite cross measure proves
holomorphy.  No value or continuity at exponent zero is asserted.
\end{proof}
